\documentclass{template}

\usepackage{mathtools}
\usepackage{mdframed}
\usepackage{varwidth} 
\pgfplotsset{compat=1.18}

\newcommand{\qedwhite}{\hfill \ensuremath{\Box}}

% Random sampling
\newcommand\leftarrowS{\leftarrow\joinrel\smalldollar}
\newcommand\rightarrowS{\smalldollar\joinrel\rightarrow}

\makeatletter
\newcommand{\smalldollar}{\mathrel{\mathpalette\small@dollar\relax}}
\newcommand{\small@dollar}[2]{%
  \vcenter{\hbox{%
    $#1\textnormal{\fontsize{0.7\dimexpr\f@size pt}{0}\selectfont\$}$%
  }}%
}
\makeatother


\begin{document}

\begin{center}
    \Large\textcolor{blue!70!black}{\textbf{Introduction to Lattice-based Cryptography}} \\
    \large \textbf{Mathematics Seminar at Kyiv School of Economics} \\
    \normalsize\textit{Dmytro Zakharov}
\end{center}

\begin{abstract}
    In recent years, the interest in lattice-based cryptography has been
    steadily increasing. One primary reason is, of course, that problems related
    to lattices and their derivatives are widely believed to be \textit{hard
    even for quantum computers}. Furthermore, there are settings where the
    algebraic structure of lattices is natural: e.g., for fully-homomorphic
    encryption.

    This brief seminar is devoted to technical introduction to lattice-based
    cryptography. We give formal definitions of lattices and introduce key
    related hardness assumptions such as SVP, CVP, and LWE. We show how to build
    public-key encryption (PKE) scheme and give an overview of basic ideas of
    modern cryptanalysis.
\end{abstract}

\tableofcontents

\newpage

\section{Lattices}\label{section:lattices}

\subsection*{Notation and Basic Notions}

All vectors are written in bold (e.g., $\mathbf{x}$ or $\mathbf{y}$) and
considered \textit{column} vectors by default. All real vector spaces are
Euclidean: that is, for every $\mathbf{x}, \mathbf{y} \in V \subseteq
\mathbb{R}^d$, we define the inner product $\langle \cdot, \cdot \rangle$ and
$\ell_2$ norm $\|\cdot\|$ in a standard fashion:
\begin{equation*}
    \langle \mathbf{x}, \mathbf{y} \rangle \triangleq \sum_{i \in [d]} x_iy_i, \quad \|\mathbf{x}\| \triangleq \sqrt{\langle \mathbf{x}, \mathbf{x} \rangle},
\end{equation*}

where we use $[d]$ to denote the set $\{1,\dots,d\}$. Additionally, we use $(d)$
to denote the ranging set using 0-indexed notation: $\{0,\dots,d-1\}$. Other
norms will be introduced on the go when needed. Finally, we use notation $x
\leftarrowS \mathcal{X}$ to denote sampling from the distribution $\mathcal{X}$.
If $\mathcal{X}$ is a set, the underlying distribution is uniform over
$\mathcal{X}$. 

\subsection{Definition}

We start our discussion with the definition of the lattice. Following
\cite{nist}, we provide two definitions: one is rather algebraic, while the
other is more intuitive and practical. In fact, both definitions will be helpful
since some facts are sometimes easier to prove using the first definition. 
\begin{definition}\label{def:lattice}
    Let $V$ be a real Euclidean vector space of finite dimension $d$. The
    following two definitions of \textbf{lattice} $\Lambda \subseteq V$ of rank
    $r \leq d$ are equivalent:
    \begin{enumerate}
        \item[\textit{(i)}] $\Lambda$ is a discrete additive subgroup of $V$
        (with $\Lambda \cong \mathbb{Z}^r$ as $\mathbb{Z}$-modules).
        \item[\textit{(ii)}] $\Lambda = \mathcal{L}(B) \triangleq \{B\mathbf{z}:
        \mathbf{z} \in \mathbb{Z}^r\}$ for some matrix $B \in \mathbb{R}^{d
        \times r}$ where its columns, further denoted as
        $\{\mathbf{b}_1,\dots,\mathbf{b}_r\} \subseteq \mathbb{R}^d$, are
        linearly independent over $V$. 
    \end{enumerate}
\end{definition}

While the second definition is clear, the first one requires some deciphering.
Since we work in the Euclidean space, it makes sense to define what it means for
the set $S \subseteq V$ to be \textit{discrete}. Intuitively, it should mean
that any two points $\mathbf{s}_1,\mathbf{s}_2 \in S$ are separated by some
``considerable'' distance. To translate it into the mathematical language, we
require that:
\begin{equation*}
    \inf_{\mathbf{s}_1 \neq \mathbf{s}_2 \in S} \|\mathbf{s}_1-\mathbf{s}_2\| > 0.
\end{equation*}

Besides $\Lambda$ being discrete, additive subgroup structure of $\Lambda$
guarantees that $\Lambda$ is closed under any $\mathbb{Z}$-linear combinations:
that is, for any $\mathbf{x}, \mathbf{y} \in \Lambda$ and $\lambda, \mu \in
\mathbb{Z}$, we have $\lambda\mathbf{x} + \mu\mathbf{y} \in \Lambda$. Showing
this property might sometimes be easier than finding the matrix $B$ in
definition~\textit{(ii)} explicitly.

\textbf{Proof of definition equivalence.} $(ii) \Rightarrow (i)$. Clearly we
have an additive subgroup structure. Discreteness comes from the fact that since
$B$ has full column rank, there exists some $\gamma$ such that $\|B\mathbf{z}\|
\geq \gamma\|\mathbf{z}\|$ (in fact, $\gamma$ is the smallest singular value)
for all $\mathbf{z} \in \mathbb{R}^r$. Thus, for $\mathbf{z} \in \mathbb{Z}^r
\setminus \{0\}$, one has $\|B\mathbf{z}\| \geq \gamma > 0$. 

$(i) \Rightarrow (ii)$. Assume $\Lambda \cong \mathbb{Z}^r$ is a discrete
additive subgroup of $V$. Since $\Lambda$ is free abelian of rank $r$, there
exist elements $\mathbf{b}_1,\dots,\mathbf{b}_r \in \Lambda$ such that any
$\lambda \in \Lambda$ can be written as $\lambda = \sum_{i \in
[r]}z_i\mathbf{b}_i$, $z_i \in \mathbb{Z}$. Now we shall prove linear
independence of $(\mathbf{b}_1,\dots,\mathbf{b}_r)$ over $\mathbb{R}$. 
In that is the case, then $\Lambda = \mathcal{L}(B)$ and we are done.

Assume otherwise that there is a non-trivial linear combination $\sum_{i \in
[r]}\alpha_i\mathbf{b}_i=0$. Pick $\{q_i^{(n)}\}_{n \geq 0}$ that approximate
corresponding $\alpha_i$ ($\lim_{n \to \infty}q_i^{(n)} = \alpha_i$) and pick
$N_n$ be the product of all denominators of $\{q_i^{(n)}\}_{i \in [r]}$. Then,
define $\lambda_n := \sum_{i \in [r]}N_nq_i^{(n)}\mathbf{b}_i \in \Lambda$. Now
note that $\lim_{n \to \infty}\lambda_n = 0$, contradicting the discreteness. 

\textbf{Geometric Intuition.} While definition~\textit{(i)} is useful, the
definition~\textit{(ii)} gives an intuitive way of thinking of lattice $\Lambda
= \mathcal{L}(B)$. If vector space $V$ is typically thought as a set of all
$\mathbb{R}$-linear combinations\footnote{We further use notation $S+S':=\{s+s':
s \in S, s' \in S'\}$ and for $a \in S$, set $aS := \{as: s \in S\}$.} $\sum_{i
\in [d]}\mathbf{e}_i\mathbb{R}$ over some basis $\{\mathbf{e}_i\}_{i \in [d]}$,
the lattice $\Lambda$ can be thought as all $\mathbb{Z}$-linear combinations
$\sum_{i \in [r]}\mathbf{b}_i\mathbb{Z}$ over real basis $\{\mathbf{b}_i\}_{i
\in [r]} \subseteq \mathbb{R}^d$. See \Cref{fig:lattice} for illustration.

\begin{remark}
    In this section we consider \textit{full-rank} lattices only: that is,
    $\Lambda \cong \mathbb{Z}^{d}$ for $d=\dim(V)$.
\end{remark}

\begin{figure}[ht]
  \centering
  \begin{tikzpicture}
    \coordinate (Origin)   at (0,0);
    \coordinate (XAxisMin) at (-2,0);
    \coordinate (XAxisMax) at (5,0);
    \coordinate (YAxisMin) at (0,-2);
    \coordinate (YAxisMax) at (0,5);
    \draw [ultra thick, gray,-stealth] (XAxisMin) -- (XAxisMax);
    \draw [ultra thick, gray,-stealth] (YAxisMin) -- (YAxisMax);

    \clip (-2,-2) rectangle (5cm,5cm); 
    \pgftransformcm{1}{0.6}{0.7}{1}{\pgfpoint{0cm}{0cm}}

    \coordinate (Bone) at (0,2);
    \coordinate (Btwo) at (2,-2);
    \draw[style=help lines, dashed, ultra thick, black!20] (-14,-14) grid[step=2cm] (14,14);

    \foreach \x in {-7,-6,...,7}{
      \foreach \y in {-7,-6,...,7}{
        \node[draw,circle,inner sep=2pt,fill] at (2*\x,2*\y) {};
      }
    }

    \draw [ultra thick,-latex,blue!70!black] (Origin)
        -- (Bone) node [right] {$\mathbf{b}_1$};
    \draw [ultra thick,-latex,blue!70!black] (Origin)
        -- (Btwo) node [below right] {$\mathbf{b}_2$};
    \filldraw[fill=red!70!purple, dashed, very thick, fill opacity=0.3, draw=black] 
        (Origin) -- (Bone) -- ($(Bone) + (Bone) - (Btwo)$) -- ($(Bone)-(Btwo)$) -- cycle;
  
    %\draw [ultra thick,-latex,red!70!purple] (Origin)
    %    -- ($0+1*(Bone)-3*(Btwo)$) node [above left] {$\mathbf{b}_1-3\mathbf{b}_2$};
    \draw [ultra thick,-latex,red!70!purple] (Origin)
        -- ($(Bone)-(Btwo)$) node [above left] {$\mathbf{b}_1-\mathbf{b}_2$};
    \filldraw[fill=blue!70!black, dashed, very thick, fill opacity=0.3, draw=black] 
    (Origin) -- (Bone) -- ($(Bone)+(Btwo)$) -- (Btwo) -- cycle;
    \node at ($0.5*(Bone)+0.5*(Btwo)$) {\footnotesize$\mathcal{P}(\textcolor{blue!70!black}{B})$};
    \node at ($0.5*(Bone) + 0.5*($(Bone)-(Btwo)$)$) {\footnotesize$\;\mathcal{P}(\textcolor{red!70!purple}{B'})$};
  \end{tikzpicture}
  \caption{Lattice $\Lambda=\mathcal{L}(\textcolor{blue!70!black}{B})$ of rank
  $2$ with basis $\textcolor{blue!70!black}{B} =
  \{\textcolor{blue!70!black}{\mathbf{b}_1},\textcolor{blue!70!black}{\mathbf{b}_2}\}$.
  $\textcolor{blue!70!black}{B}$ is an example of a ``good'' basis where basis
  vectors are almost orthogonal. The basis $\textcolor{red!70!purple}{B'} =
  \{\textcolor{blue!70!black}{\mathbf{b}_1},\textcolor{red!70!purple}{\mathbf{b}_1-\mathbf{b}_2}\}$
  of the same lattice $\Lambda=\mathcal{L}(\textcolor{red!70!purple}{B'})$ is ``worse'' (the vectors seem more parallel and
  fundamental domain $\mathcal{P}(\textcolor{red!70!purple}{B'})$ is skewed).}
  \label{fig:lattice}
\end{figure}

\subsection{Fundamental Measures}

\subsubsection{Fundamental Domain}

A basis $B$ of $\Lambda = \mathcal{L}(B)$ is associated with the fundamental
domain $\mathcal{P}(B) \triangleq \sum_{i \in
[d]}\mathbf{b}_i[0,1)$\footnote{Sometimes, the interval $[-1/2, 1/2)$ is used
instead of $[0,1)$ to center $\mathcal{P}(B)$ with respect to origin: in such
case, $\mathcal{P}(B)$ is called a \textit{centered parallelepiped of $B$}}.

What is special about this parallelepiped is that it represents all elements in
the quotient group $V/\Lambda$: any element $\mathbf{x} \in V$ can be uniquely
written as $\mathbf{x} = \mathbf{v} + \mathbf{p}$ where $\mathbf{v} \in \Lambda$
and $\mathbf{p} \in \mathcal{P}(B)$. Moreover, the shape of this
$\mathcal{P}(B)$ determines the ``quality'' of the basis $B$: if
$\mathcal{P}(B)$ is skewed and is not concentrated in the origin, it is
typically considered ``bad'' as some computational problems become much harder
to solve. In turn, a non-skewed basis $B'$ primarily concentrated around the
origin is considered ``good'' as $\{\mathbf{b}_i'\}_{i \in [d]}$ much more
closely resembles the orthogonal basis than one of $B$ (and of course, the
orthogonal basis in a sense is ``perfect'' because it is very ``nice'' to work
with). Again, see \Cref{fig:lattice} for illustration.

Another natural quantity to consider is the volume of such parallelepiped,
called \textit{covolume}.

\begin{definition}[Covolume]
    The \textbf{covolume} $\text{Vol}(\Lambda)$ of lattice $\Lambda$ is defined as:
    \begin{equation*}
        \text{Vol}(\Lambda) \triangleq \text{Vol}(V/\Lambda) = \sqrt{|\det(B^{\top}B)|}.
    \end{equation*}
\end{definition}

Note that the definition of $\mathcal{P}(B)$ depends on the choice of $B$ since
the basis of $\Lambda$ is generally not unique. Thus, it might seem that
covolume should also depend on the choice of $B$. However, we show the
following.

\begin{proposition}
    The \textbf{covolume} $\text{Vol}(\Lambda)$ is an invariant of $\Lambda$:
    that is, it does not depend on the choice of $B$.
\end{proposition}

\textbf{Proof.} Suppose $\Lambda = \mathcal{L}(B_0) = \mathcal{L}(B_1)$. Since
$B_0$ and $B_1$ both generate $\Lambda$, one has $B_1 = UB_0$ for some matrix $U
\in \text{GL}_d(\mathbb{Z})$. Since $U$ is an integer matrix with integer
inverse, it is unimodular: that is, $|\det U| = 1$. It is then easy to see
$|\det B_1^{\top}B_1| = |\det B_0^{\top}U^{\top}UB_0| = |\det B_0^{\top}B_0| =
\text{Vol}(\Lambda)^2$. \qedwhite

\subsubsection{Successive Minima}

Turns out that one of the hardest computational problems is computing the
shortest vectors of $\Lambda$. However, we need some preliminary work to
formalize this statement. 

Denote $B_{\rho}(\mathbf{x}_0) := \{\mathbf{x} \in V:
\|\mathbf{x}-\mathbf{x}_0\|<\rho\}$ --- an open ball in $V$ of radius $\rho$
centered at $\mathbf{x}_0$. 

\begin{definition}[Successive minima]
    The \textbf{$i^{\text{th}}$ minimum $\lambda_i(\Lambda)$} is the radius 
    of the smallest sphere centered at the origin containing $i$ linearly 
    independent lattice vectors:
    \begin{equation*}
        \lambda_i(\Lambda) \triangleq \inf \{\rho \in \mathbb{R}_{\geq 0}: \dim(\text{span}(\Lambda \cap B_{\rho}(\mathbf{0}))) \geq i\}.
    \end{equation*}
\end{definition}

In particular, $\lambda_1(\Lambda) = \min_{\mathbf{x} \neq \mathbf{y} \in
\Lambda}\|\mathbf{x}-\mathbf{y}\| = \min_{\mathbf{x} \in
\Lambda\setminus\{\mathbf{0}\}}\|\mathbf{x}\|$ is called the \textit{first
minimum} while the vector of length $\lambda_1(\Lambda)$ --- the shortest
vector. The natural question to ask though is whether it exists at all: after
all, in $\mathbb{R}^d$, for instance, it is clearly undefined. We give
affirmative answer.

We shall further denote the result of Gram-Schmidt orthogonalization of $B$ by
$B^*$. The question that might naturally arise: given such orthogonalization,
can't we then always turn a ``bad'' basis of lattice $\Lambda$ into a ``good''
basis? Unfortunately, typically $\mathcal{L}(B) \neq \mathcal{L}(B^*)$: the
orthogonalization of a lattice basis is (typically) not a lattice basis.
However, the idea of Gram-Schmidt orthogonalization lies in the core of many
attacks on lattices. 
\begin{proposition}\label{prop:gram-schmidt-props} Suppose $\Lambda =
    \mathcal{L}(B)$ is a lattice and $B^*$ is the Gram-Schmidt orthogonalization
    of $B$ with columns $\{\mathbf{b}^*_j\}_{j \in [d]}$. Then
    $\lambda_1(\Lambda) \geq \min_{j \in [d]}\|\mathbf{b}_j^*\| > 0$.
\end{proposition}

\textbf{Proof} (adapted from~\cite{complexity-goldwasser}). Consider a generic
non-zero lattice vector $B\mathbf{x}$ ($\mathbf{x} \in \mathbb{Z}^d \setminus
\{\mathbf{0}\}$) and let $i$ be the biggest index such that $x_i \neq 0$. We
first prove that $|\langle B\mathbf{x}, \mathbf{b}_i^* \rangle| \geq
\|\mathbf{b}_i^*\|^2$. From the definition of $i$, we know $B\mathbf{x} =
\sum_{j=1}^i \mathbf{b}_jx_j$. Thus:
\begin{align*}
    \langle B\mathbf{x}, \mathbf{b}_i^* \rangle = \sum_{j=1}^i \langle \mathbf{b}_j, \mathbf{b}_i^*\rangle x_j &= \langle \mathbf{b}_i,\mathbf{b}_i^* \rangle x_i\\
    &= \langle \mathbf{b}_i^* + \sum_{j<i}\mu_{i,j}\mathbf{b}_j^*, \mathbf{b}_i^*\rangle x_i \\
    &= \langle \mathbf{b}_i^*, \mathbf{b}_i^*\rangle x_i + \sum_{j < i}\mu_{i,j}\langle \mathbf{b}_j^*, \mathbf{b}_i^*\rangle x_i \\
    &= \|\mathbf{b}_i^*\|^2x_i.
\end{align*}
From this follows $|\langle B\mathbf{x}, \mathbf{b}_i^* \rangle| \geq
\|\mathbf{b}_i^*\|^2|x_i| \geq \|\mathbf{b}_i^*\|^2$ and therefore
$\|B\mathbf{x}\|\cdot\|\mathbf{b}_i^*\| \geq \|\mathbf{b}_i^*\|^2$. Clearly
$\|\mathbf{b}_i^*\| \neq 0$ since $\mathbf{b}_i$'s are linearly independent, so
$\|B\mathbf{x}\| \geq \|\mathbf{b}_i^*\| \geq \min_{j \in
[d]}\|\mathbf{b}_j^*\|$. By definition of $\lambda_1(\Lambda)$, the result
follows. \qedwhite

\begin{proposition}
    For lattice $\Lambda$, there exists $\mathbf{v} \in \Lambda$ such that
    $\|\mathbf{v}\| = \lambda_1(\Lambda)$.
\end{proposition}

\textbf{Proof.} By the definition of $\lambda_1(\Lambda)$, there exists a
sequence $\{\mathbf{v}_n\}_{n \in \mathbb{N}} \subseteq \Lambda$ such that
$\lim_{n \to \infty}\|\mathbf{v}_n\| = \lambda_1(\Lambda)$. From
Proposition~\ref{prop:gram-schmidt-props}, $\lambda_1(\Lambda) > 0$, so for all
sufficiently large $n$ it must be $\|\mathbf{v}_n\| \leq 2\lambda_1(\Lambda)$,
i.e., lattice vector $\mathbf{v}_n$ belongs to the closed ball
$\overline{B}_{2\lambda_1(\Lambda)}(\mathbf{0})$. But this set is compact, so we
might extract a convergent subsequence $\{\mathbf{v}_{i_j}\}_{j \in \mathbb{N}}$
with limit $\mathbf{w} = \lim_{j \to \infty}\mathbf{v}_{i_j}$ and
$\|\mathbf{w}\| = \lambda_1(\Lambda)$. We claim that $\mathbf{w} \in \Lambda$.
By definition of $\mathbf{w}$, $\lim_{j \to \infty}\|\mathbf{v}_{i_j} -
\mathbf{w}\|=0$. Therefore for all sufficiently large $j$, $\|\mathbf{v}_{i_j} -
\mathbf{w}\| < \lambda_1(\Lambda)/2$. By triangle inequality for sufficiently
large $j$ and all $k>j$, $\|\mathbf{v}_{i_j}-\mathbf{v}_{i_k}\| \leq
\|\mathbf{v}_{i_j} - \mathbf{w}\| + \|\mathbf{w} - \mathbf{v}_{i_k}\| <
\lambda_1(\Lambda)$. But $\mathbf{v}_{i_j}-\mathbf{v}_{i_k} \in \Lambda$, and 
no non-zero lattice vector can have length strictly less than $\lambda_1(\Lambda)$.
Thus, $\mathbf{v}_{i_k}=\mathbf{v}_{i_j}$ for all $k>j$. Thus $\mathbf{w}=\mathbf{v}_{i_j} \in \Lambda$. \qedwhite

\subsection{Minkowski's Theorem}

Clearly, one expects the dependence between the covolume of a lattice
$\text{Vol}(\Lambda)$ and the first minimum $\lambda_1(\Lambda)$. In particular,
for smaller covolume we expect the smaller first minimum $\lambda_1(\Lambda)$
(or alternatively, one can say the lattice with large covolume cannot have a
small first minimum). Moreover, by dimensionality analysis one should expect
that $\lambda_1(\Lambda) \propto \text{Vol}(\Lambda)^{1/d}$. The Minkowski
theorem gives the inequality between these two quantities.

\begin{theorem}[Minkowski's Theorems]\label{theorem:minkowski}
    Let $\Lambda=\mathcal{L}(B)$ be the lattice of rank $d$ with covolume
    $\text{Vol}(\Lambda) = \det B$. Then,
    \begin{enumerate}[align=left]
        \item[\textbf{First Minkowski's Theorem}] $\lambda_1(\Lambda) < \sqrt{d}\text{Vol}(\Lambda)^{1/d}$.
        \item[\textbf{Second Minkowski's Theorem}] $\prod_{i \in [d]}\lambda_i^{1/d} < \sqrt{d}\text{Vol}(\Lambda)^{1/d}$. 
    \end{enumerate}
\end{theorem}

\textbf{Proof.} See \cite[Section 1.3]{complexity-goldwasser}.

While the Minkowski's Theorem turns out to be asymptotically tight for ``bad''
lattices, in general $\lambda_1(\Lambda)$ might be much smaller than
$\sqrt{d}\text{Vol}(\Lambda)^{1/d}$. Consider a very simple example: let $d=2$
and $\mathbf{b}_1 := \varepsilon \mathbf{e}_1$, $\mathbf{b}_2 :=
(1/\varepsilon)\mathbf{e}_2$ (where $\mathbf{e}_1,\mathbf{e}_2$ is the canonical
basis for $\mathbb{R}^2$). Minkowski's Theorem gives $\lambda_1(\Lambda) \leq
\sqrt{2}$, however one easily sees that $\lambda_1(\Lambda) = \varepsilon$ can
be made arbitrarily small.

Moreover, Minkowski's Theorem proof is not constructive: we know that a vector
of length below $\sqrt{d}\text{Vol}(\Lambda)^{1/d}$ exists, but there is no
concrete algorithm to find it, let alone $\lambda_1(\Lambda)$. This became the
core of one of the most well-known lattice computational problem: finding
shortest vector in $\Lambda$. We introduce this and related problems in the next
subsection.

\subsection{Lattice Computational Problems}

In this section, we finally define some problems that are currently believed to
be unsolvable in polynomial time (under reasonable choice of parameters, of
course). 

\begin{itemize}
    \item Finding $\mathbf{v}$ such that $\|\mathbf{v}\| = \lambda_1(\Lambda)$
    of $\Lambda = \mathcal{L}(B)$ from a random basis $B$.
    \item Finding $\mathbf{v}$ such that $\|\mathbf{v}\| \leq
    \gamma\lambda_1(\Lambda)$.
    \item Given $\mathbf{t} \leftarrowS \mathbb{Z}^{\mu}$, find some $\mathbf{v}
    \in \mathcal{L}$ such that $\text{dist}(\mathbf{t}, \mathbf{v}) \leq \gamma
    \cdot \text{dist}(\mathbf{t}, \mathcal{L})$.
\end{itemize}

\section{Learning With Errors (LWE)}\label{section:lwe}

\subsection{Key Generation}

Fix some modulo $q>1$ (further assume it is prime). Our primary goal is to build
public key cryptography schemes in the lattice-fashionable way. Turns out that
one of such approaches is to use \textbf{Learning With Errors} assumptions.

Fix secret key $\mathbf{s} \in \mathbb{Z}_q^n$ which we assume, for the moment,
is sampled uniformly randomly from the corresponding set. The public key,
corresponding to $\mathbf{s}$, will be a vector of perturbed random inner
products with $\mathbf{s}$. Specifically, the public information consists of
multiple pairs $(\mathbf{a}, b) \in \mathbb{Z}_q^{n} \times \mathbb{Z}_q$ such
that $b = \langle\mathbf{a},\mathbf{s}\rangle + e$ for an \textit{error} $e \in
\mathbb{Z}_q$, typically taken from the distribution that produces very small
values of $e$: for instance, one might think of $e \in [-\eta,\eta]\pmod{q}$ for
$\eta \approx 3$. Denote such distribution by $\chi$.

This way, we can give an intermediate version of our key generation as follows:

\begin{center}
\begin{varwidth}{0.7\linewidth}
\begin{mdframed}[linecolor=black,
                 linewidth=1.0pt,
                 skipabove=10pt,
                 skipbelow=10pt,
                 leftmargin=0pt,
                 rightmargin=0pt]
\colorbox{purple!10}{\underline{\textbf{$\mathsf{KeyGen}(1^{\lambda}) \to (\mathsf{pk},\mathsf{sk})$:}}}

\begin{enumerate}
    \item $\mathbf{s} \leftarrowS \mathbb{Z}_q^{\lambda}$; \hspace{5px}{\textcolor{gray!70!black}{\footnotesize\texttt{// Generate secret}}}
    \item \underline{\textbf{for}} $i \in [\lambda]$:
    \item \hspace{10px} $\mathbf{a}_i \leftarrowS \mathbb{Z}_q^{\lambda}$, $e \leftarrowS \chi$;
    \item \hspace{10px} $b_i \gets \langle\mathbf{a}_i,\mathbf{s}\rangle + e$;\hspace{5px}{\footnotesize\texttt{\textcolor{gray!70!black}{// Generate random pertubation}}}
    \item \textbf{return} $(\mathsf{pk}=\{(\mathbf{a}_i, b_i)\}_{i \in [\lambda]}, \mathsf{sk}=\mathbf{s})$.
\end{enumerate}
\end{mdframed}
\end{varwidth}
\end{center}

It is then natural to batch all pertubations into a single matrix $A$ and vector
$\mathbf{b}$.\footnote{We also note that $A$ might not necessarily be (and is
typically not) the square matrix, but just for the sake of exposition we for now
assume $A \in \mathbb{Z}_q^{\lambda \times \lambda}$.} This way, our updated key
generation scheme is in \Cref{fig:keygen} (where by $\chi^{\lambda}$ we denote
the distribution over $\mathbb{Z}_q^{\lambda}$ where each component is taken
from $\chi$).

\begin{figure}[ht]
\centering
\begin{varwidth}{0.4\linewidth}
\begin{mdframed}[linecolor=black,
                 linewidth=1.0pt,
                 skipabove=10pt,
                 skipbelow=10pt,
                 leftmargin=0pt,
                 rightmargin=0pt]
\colorbox{purple!10}{\underline{$\mathsf{KeyGen}(1^{\lambda}) \to (\mathsf{pk},\mathsf{sk})$:}}
\begin{enumerate}
    \item $(\mathbf{s}, A) \leftarrowS \mathbb{Z}_q^{\lambda} \times \mathbb{Z}_q^{\lambda \times \lambda}$;
    \item $\mathbf{e} \leftarrowS \chi^{\lambda}$;
    \item $b \gets A\mathbf{s} + \mathbf{e}$;
    \item \textbf{return} $(\mathsf{pk}=(A,b), \mathsf{sk}=\mathbf{s})$.
\end{enumerate}
\end{mdframed}
\end{varwidth}
\caption{LWE-based Key Generation Scheme}
\label{fig:keygen}
\end{figure}

Note that we of course expect that solving $A\mathbf{s}+\mathbf{e}=\mathbf{b}$
is hard with respect to $\mathbf{s}$. It should also be clear why we introduce
the error term $\mathbf{e}$: assuming $\mathbf{e}=\mathbf{0}$, the equation
simplifies to $A\mathbf{s}=\mathbf{b}$. This is a simple system of linear
equations which can be further dealt with by means of, say, Gaussian
elimination, which has $O(\lambda^3)$ complexity. The assumption of hardness of
solving this equation for $\mathbf{e} \neq \mathbf{0}$ is what we introduce and
treat in the next sections.

\subsection{Search-LWE (SLWE) and Decisional-LWE (DLWE) Assumptions}

In this section we shall finally formalize main assumptions about learning with
errors problem. First, introduce the LWE instance generation algorithm as
specified in~\Cref{fig:lwe-instance} below.

\begin{figure}[ht]
\centering
\begin{varwidth}{0.45\linewidth}
\begin{mdframed}[linecolor=black,
                 linewidth=1.0pt,
                 skipabove=10pt,
                 skipbelow=10pt,
                 leftmargin=0pt,
                 rightmargin=0pt]
\colorbox{purple!10}{\underline{$\mathsf{LWEGen}_{n,m,q,\chi}(1^{\lambda}, \mathbf{s}) \to \mathbb{Z}_q^{m \times n} \times \mathbb{Z}_q^m$:}}
\begin{enumerate}
    \item $A \leftarrowS \mathbb{Z}_q^{m \times n}$;
    \item $\mathbf{e} \leftarrowS \chi$ (over $\mathbb{Z}_q^m$);
    \item $\mathbf{b} \gets A\mathbf{s} + \mathbf{e}$;
    \item \textbf{return} $(A, \mathbf{b})$.
\end{enumerate}
\end{mdframed}
\end{varwidth}
\caption{LWE Instance Generation}
\label{fig:lwe-instance}
\end{figure}

Now, informally, LWE assumption states that, given output of
$\mathsf{LWEGen}_{n,m,q,\chi}(1^{\lambda},\mathbf{s})$, it is impossible to
determine $\mathbf{s}$ (\textit{Search-LWE assumption}). Similarly to many
assumptions in other cryptographic systems (such as Diffie-Hellman assumption),
LWE also has the decisional analog: distinguish output of
$\mathsf{LWEGen}_{n,m,q,\chi}(1^{\lambda},\mathbf{s})$ from the uniform sample
over $\mathbb{Z}_q^{m \times n} \times \mathbb{Z}_q^m$ (\textit{Decisional-LWE
assumption}).  We state both assumptions more formally below.

\begin{definition}[Search-LWE Assumption]\label{def:search-lwe-assumption}
    For the given adversary $\mathcal{A}$ and LWE parameters $(n,m,q,\chi)$,
    define $\underline{\mathsf{Game} \;
    \mathsf{SLWE}}_{\mathsf{LWEGen}}^{\mathcal{A}}$ as specified
    in~\Cref{fig:search-lwe}. Define the advantage of $\mathcal{A}$ as
    \begin{equation*}
        \mathsf{Adv}^{\text{SLWE}}_{\mathcal{A},\mathsf{LWEGen}}(\lambda) = \text{Pr}\left[\underline{\mathsf{Game} \;
    \mathsf{SLWE}}_{\mathsf{LWEGen}}^{\mathcal{A}}(\lambda) = 1\right].
    \end{equation*}

    We say that \textbf{Search-LWE assumption} (SLWE) holds if
    $\mathsf{Adv}^{\text{SLWE}}_{\mathcal{A},\mathsf{LWEGen}}(\lambda)=\text{negl}(\lambda)$.
\end{definition}

\begin{definition}[Decisional-LWE Assumption]\label{def:decision-lwe-assumption}
    For the given adversary $\mathcal{A}$ and LWE parameters $(n,m,q,\chi)$,
    define $\underline{\mathsf{Game} \;
    \mathsf{DLWE}}_{\mathsf{LWEGen}}^{\mathcal{A}}$ as specified
    in~\Cref{fig:decision-lwe}. Define the advantage of $\mathcal{A}$ as
    \begin{equation*}
        \mathsf{Adv}^{\text{DLWE}}_{\mathcal{A},\mathsf{LWEGen}}(\lambda) = \left|\text{Pr}\left[\underline{\mathsf{Game} \;
    \mathsf{DLWE}}_{\mathsf{LWEGen}}^{\mathcal{A}}(\lambda) = 1\right] - \frac{1}{2}\right|.
    \end{equation*}

    We say that \textbf{Decisional-LWE assumption} (DLWE) holds if
    $\mathsf{Adv}^{\text{DLWE}}_{\mathcal{A},\mathsf{LWEGen}}(\lambda)=\text{negl}(\lambda)$
    for all PPT adversaries $\mathcal{A}$.
\end{definition}

\begin{figure}[ht]
    \centering
    \begin{minipage}{0.45\linewidth}
        \begin{mdframed}[linecolor=black, linewidth=1pt,
                         skipabove=10pt, skipbelow=10pt]
        \colorbox{green!10}{\underline{$\mathsf{Game} \; \mathsf{SLWE}_{\mathsf{LWEGen}}^{\mathcal{A}}$:}}
        \begin{enumerate}
            \item $\mathbf{s} \leftarrowS \mathbb{Z}_q^{n}$;
            \item $(A,\mathbf{b}) \leftarrowS \mathsf{LWEGen}_{n,m,q,\chi}(\mathbf{s})$;
            \item $\widehat{\mathbf{s}} \gets \mathcal{A}(A,\mathbf{b})$;
            \item \textbf{return} $\mathbf{s} =_? \widehat{\mathbf{s}}$.
        \end{enumerate}
        \end{mdframed}
        \caption{Search-LWE Game} 
        \label{fig:search-lwe}
    \end{minipage}
    \hfill
    \begin{minipage}{0.45\linewidth}
        \begin{mdframed}[linecolor=black, linewidth=1pt,
                         skipabove=10pt, skipbelow=10pt]
        \colorbox{green!10}{\underline{$\mathsf{Game} \; \mathsf{DLWE}_{\mathsf{LWEGen}}^{\mathcal{A}}$:}}
        \begin{enumerate}
            \item $\mathbf{s} \leftarrowS \mathbb{Z}_q^{n}$;
            \item $(A_0,\mathbf{b}_0) \leftarrowS \mathsf{LWEGen}_{n,m,q,\chi}(\mathbf{s})$;
            \item $(A_1,\mathbf{b}_1) \leftarrowS \mathbb{Z}_q^{m \times n} \times \mathbb{Z}_q^m$;
            \item $c \leftarrowS \{0,1\}$;
            \item $\widehat{c} \gets \mathcal{A}(A_c,\mathbf{b}_c)$;
            \item \textbf{return} $c =_? \widehat{c}$.
        \end{enumerate}
        \end{mdframed}
        \caption{Decision-LWE Game}
        \label{fig:decision-lwe}
    \end{minipage}
\end{figure}

The natural question to ask is whether SLWE and DLWE assumptions are equivalent
or not. In fact, \cite{regev} has proven that these notions are in fact
equivalent in case $q=\mathsf{poly}(n)$. We show this in the following
proposition.

\begin{proposition}[Following~\cite{regev}]
    If $q=\mathsf{poly}(n)$, SLWE and DLWE assumptions are polynomially
    equivalent. To put it the other way, if the SLWE assumption holds for LWE
    instance with parameters $(n,m,q=\mathsf{poly}(n),\chi)$, then the DLWE
    assumption also holds.
\end{proposition}

\textbf{Proof.} $\text{SLWE} \Rightarrow \text{DLWE}$ reduction is obvious. We
therefore show $\text{DLWE} \Rightarrow \text{SLWE}$. Suppose exists adversary
$\mathcal{A}$ that wins $\underline{\mathsf{Game} \;
\mathsf{DLWE}}_{\mathsf{LWEGen}}^{\mathcal{A}}$ in a non-negligible probability.
We construct $\mathcal{B}^{\mathcal{A}}$ that wins $\underline{\mathsf{Game} \;
\mathsf{SLWE}}_{\mathsf{LWEGen}}^{\mathcal{B}}$ is a non-negligible probability
as well.

The informal idea is as follows: suppose $\mathcal{B}$ wants to solve
$A\mathbf{s}+\mathbf{e}=\mathbf{b}$. Turns out that $\mathcal{B}$ can guess the
coordinates of $\mathbf{s}$ one at a time, starting from $s_1$. Suppose
$\mathcal{B}$ picks some $\widetilde{s} \in \mathbb{Z}_q$ and wants to check
whether $\widetilde{s} = s_1$. Select some $\boldsymbol{\Delta} \leftarrowS
\mathbb{Z}_q^m$ and let
$\widetilde{A}=(\widetilde{\mathbf{a}}_1,\mathbf{a}_2,\dots,\mathbf{a}_n)$ where
$\widetilde{\mathbf{a}}_1=\mathbf{a}_1+\boldsymbol{\Delta}$ (and where the
initial matrix is given by $A=(\mathbf{a}_1,\dots,\mathbf{a}_n)$) and set
$\widetilde{\mathbf{b}} := \mathbf{b} + \widetilde{s}\boldsymbol{\Delta}$. Now,
if $\widetilde{s}=s_1$, then $\widetilde{\mathbf{b}} = \widetilde{A}\mathbf{s} +
\mathbf{e}$, otherwise $\widetilde{\mathbf{b}} = \widetilde{A}\mathbf{s} +
\mathbf{e}+(\widetilde{s}-s_1)\boldsymbol{\Delta}$. Since $\boldsymbol{\Delta}$
is uniform, it follows that in the latter case $\widetilde{\mathbf{b}}$ is
uniformly random and independent of $\widetilde{A}$. Thus, using $\mathcal{A}$,
adversary $\mathcal{B}$ has an effective test for chosen $\widetilde{s}$, and
thus $\mathcal{B}$ can range over all possible $q$ values, which is polynomial
in $n$.

This way, by using $\mathcal{B}^{\mathcal{A}}_{\text{SLWE}}$ defined in
\Cref{fig:dlwe-to-slwe}, we obtain:
\begin{equation*}
    \mathsf{Adv}^{\text{DLWE}}_{\mathcal{A},\mathsf{LWEGen}}(\lambda) \leq qn \cdot \mathsf{Adv}^{\text{SLWE}}_{\mathcal{B}^{\mathcal{A}},\mathsf{LWEGen}}(\lambda). \quad \qedwhite
\end{equation*}

\begin{remark}
    Note that it might seem like the factor of $qn$ is too large to state the
    equivalence of DLWE and SLWE (indeed, in the elliptic curve cryptography,
    one typically uses $q$ exponential in the security parameter). However, in practical
    applications, $q$ is typically not considerably large: for instance, in
    Kyber~\cite{kyber}, $q=7681$ and $n=256$!

    Also note that the brute-force complexity without the DLWE oracle is $q^n$,
    which is of course overwhelmingly large! (compared to $qn$).
\end{remark}

\begin{figure}[ht]
\centering
\begin{varwidth}{0.7\linewidth}
\begin{mdframed}[linecolor=black,
                 linewidth=1.0pt,
                 skipabove=10pt,
                 skipbelow=10pt,
                 leftmargin=0pt,
                 rightmargin=0pt]
\colorbox{purple!10}{\underline{$\mathcal{B}^{\mathcal{A}}_{\text{SLWE}}(1^{\lambda}, A=(\mathbf{a}_1,\dots,\mathbf{a}_n), \mathbf{b}) \to \mathbb{Z}_q^{m}$:}}
\begin{enumerate}
    \item $\widehat{\mathbf{s}} \gets \mathbf{0} \in \mathbb{Z}_q^m$;\hspace{5px}{\textcolor{gray!70!black}{\footnotesize\texttt{// Prepare the output}}}
    \item \textbf{for} $i \in [n]$:
    \item \hspace{10px} \textbf{for} $\widetilde{s} \in \mathbb{Z}_q$:
    \item \hspace{20px} $\boldsymbol{\Delta} \leftarrowS \mathbb{Z}_q^m$;\hspace{5px}{\textcolor{gray!70!black}{\footnotesize\texttt{// Generate random offset}}}
    \item \hspace{20px} $\widetilde{A} \gets (\mathbf{a}_1,\dots,\mathbf{a}_{i-1},\widetilde{\mathbf{a}}_i, \mathbf{a}_{i+1},\dots,\mathbf{a}_n)$, $\widetilde{\mathbf{a}}_i := \mathbf{a}_i+\boldsymbol{\Delta}$;
    \item \hspace{20px} $\widetilde{\mathbf{b}} \gets \mathbf{b} + \widetilde{s}\boldsymbol{\Delta}$;\hspace{5px}{\textcolor{gray!70!black}{\footnotesize\texttt{// Set offset vector $\widetilde{\mathbf{b}}$ and matrix $\widetilde{A}$}}}
    \item \hspace{20px} $c \gets \mathcal{A}(\widetilde{A}, \widetilde{\mathbf{b}})$;\hspace{5px}{\textcolor{gray!70!black}{\footnotesize\texttt{// Run DLWE oracle}}}
    \item \hspace{20px} \textbf{if} $c=0$, set $\widehat{s}_i \gets \widetilde{s}$ and \textbf{break};
    \item \textbf{return} $\widehat{\mathbf{s}}$.
\end{enumerate}
\end{mdframed}
\end{varwidth}
\caption{Algorithm for finding the solution to LWE instance given oracle to the DLWE.}
\label{fig:dlwe-to-slwe}
\end{figure}

\subsection{LWE-based Public-Key Encryption}

Finally, let us define the public-key encryption (PKE) scheme based on SLWE and
DLWE assumptions! We remind the reader of the definition of a PKE scheme.
\begin{definition}[Public-Key Encryption Scheme]
    A \textbf{public-key encryption scheme} $\Pi_{\text{PKE}}$ is a triplet 
    of the following PPT algorithms:
    \begin{itemize}
        \item $\mathsf{KeyGen}(1^{\lambda}) \to (\mathsf{pk},\mathsf{sk})$
        returns the pair of public key and secret key, respectively.
        \item $\mathsf{Enc}(\mathsf{pk}, m) \to c$ takes the public key and message
        and produces the ciphertext $c$.
        \item $\mathsf{Dec}(\mathsf{sk}, c) \to m$ takes the secret key and ciphertext and restores the message $m$.
    \end{itemize}
    We require \textbf{correctness}:
    \begin{equation*}
        \forall m: \; \text{Pr}\left[ \widehat{m} = m \;\middle|\; \begin{aligned}
            &(\mathsf{pk}, \mathsf{sk}) \gets \mathsf{KeyGen}(1^{\lambda}) \\
            &\widehat{m} \gets \mathsf{Dec}(\mathsf{sk}, \mathsf{Enc}(\mathsf{pk}, m))
        \end{aligned} \right] \geq 1 - \text{negl}(\lambda),
    \end{equation*}
    and \textbf{IND-CCA} and \textbf{IND-CPA} security. To not overwhelm the reader 
    with the details, we will only focus on the IND-CPA security with the corresponding $\underline{\mathsf{Game} \; \mathsf{INDCPA}}^{\mathcal{A}}_{\Pi}$ defined in \Cref{fig:ind-cpa-game}. We 
    naturally define advantage $\mathsf{Adv}_{\mathcal{A},\Pi_{\text{PKE}}}^{\text{IND-CPA}}(\lambda) := 2\left|\text{Pr}[\underline{\mathsf{Game} \; \mathsf{INDCPA}}^{\mathcal{A}}_{\Pi_{\text{PKE}}} = 1] - \frac{1}{2}\right|$ and say that $\Pi_{\text{PKE}}$ is \textbf{IND-CPA secure} if $\mathsf{Adv}_{\mathcal{A},\Pi_{\text{PKE}}}^{\text{IND-CPA}}(\lambda) = \text{negl}(\lambda)$ for all PPT $\mathcal{A}$.
\end{definition}

\begin{figure}[ht]
\centering
\begin{varwidth}{0.75\linewidth}
\begin{mdframed}[linecolor=black, linewidth=1pt,
                    skipabove=10pt, skipbelow=10pt]
\colorbox{green!10}{\underline{$\mathsf{Game} \; \mathsf{INDCPA}_{\Pi}^{\mathcal{A}}(1^{\lambda})$:}}
    \begin{enumerate}
        \item $(\mathsf{pk},\mathsf{sk}) \gets \mathsf{KeyGen}(1^{\lambda})$;
        \item $(m_0,m_1) \gets \mathcal{A}(\mathsf{pk})$;\hspace{5px}{\textcolor{gray!70!black}{\footnotesize\texttt{// Typically one requires $|m_0|=|m_1|$}}}
        \item $b \leftarrowS \{0,1\}$, $c_b \gets \mathsf{Enc}(\mathsf{pk}, m_b)$;
        \item $\widehat{b} \gets \mathcal{A}(\mathsf{pk}, m_0, m_1, c_b)$;
        \item \textbf{return} $\widehat{b} =_? b$.
    \end{enumerate}
\end{mdframed}
\end{varwidth}
\caption{IND-CPA Security Game. $\mathcal{A}$ might be
allowed to perform steps 2--3 $\text{poly}(\lambda)$ times.} 
\label{fig:ind-cpa-game}
\end{figure}

Now, we define the following PKE scheme based on $(n,m,q,\chi)$ assumption where
the message space is $\{0,1\}$. The main idea is that encryption scheme will
generate the new LWE instance $(A',\mathbf{b}')$ by taking a ``small'' linear
combination of publicly known $(A,\mathbf{b})$. By additionally adding
$m\lceil q/2 \rfloor$ to $\mathbf{b}'$, we obtain the encryption (further denote
by $\lceil \cdot \rfloor: \mathbb{Q} \to \mathbb{Z}$ the function that rounds a
given rational number to the nearest integer). Upon decryption, the holder of 
the secret key $\mathbf{s}$ will be able to check whether $A'\mathbf{s}-\mathbf{b}'$
produces the small error. If the term $m\lceil q/2 \rfloor$ was added to form 
$\mathbf{b}'$, the resultant difference would be large, suggesting that the original 
bit was $1$. Otherwise, decryption would output $0$. Let us put this formally in 
the definition below.

\begin{definition}[LWE-based PKE]\label{def:lwe-pke}
    Based on LWE instance with parameters $(n,m,q,\chi)$, define the PKE scheme
    $\Pi_{\text{PKE}}^{\text{LWE}}$ over message space $\{0,1\}$ as follows: $\mathsf{KeyGen}$ is specified in
    \Cref{fig:keygen}, $\mathsf{Enc}$ in \Cref{fig:lwe-enc}, and $\mathsf{Dec}$
    in \Cref{fig:lwe-dec}.
\end{definition}

Since we have to work with ``small'' and ``large'' elements of $\mathbb{Z}_q$,
by $\|x\|_{\infty}$ for $x \in \mathbb{Z}_q$ we denote $x$ if $x \leq \frac{q}{2}$, and
$q-x$ otherwise (for example, naturally, $\|q-3\|_{\infty} = 3$).

\begin{figure}[ht]
    \centering
    \begin{minipage}{0.45\linewidth}
        \begin{mdframed}[linecolor=black, linewidth=1pt,
                         skipabove=10pt, skipbelow=10pt]
        \colorbox{green!10}{\underline{$\mathsf{Enc}(\mathsf{pk}, m \in \{0,1\}) \to (c \in \mathbb{Z}_q^n \times \mathbb{Z}_q)$:}}
        \begin{enumerate}
            \item $\Delta\mathbf{r}, \Delta\mathbf{z} \leftarrowS \chi^m$, $\Delta z' \leftarrowS \chi$;
            \item $\widetilde{\mathbf{a}} \gets A^{\top}\Delta\mathbf{r} + \Delta\mathbf{z}$;
            \item $\widetilde{b} \gets \mathbf{b}^{\top}\Delta\mathbf{r}  + \Delta z' + m\lceil q/2 \rfloor$;
            \item \textbf{return} $c=(\widetilde{\mathbf{a}}, \widetilde{b})$.
        \end{enumerate}
        \end{mdframed}
        \caption{Encryption in $\Pi_{\text{PKE}}^{\text{LWE}}$} 
        \label{fig:lwe-enc}
    \end{minipage}
    \hfill
    \begin{minipage}{0.45\linewidth}
        \begin{mdframed}[linecolor=black, linewidth=1pt,
                         skipabove=10pt, skipbelow=10pt]
        \colorbox{green!10}{\underline{$\mathsf{Dec}(\mathsf{sk}, c) \to m$:}}
        \begin{enumerate}
            \item Parse $(\widetilde{\mathbf{a}}, \widetilde{b}) \gets c$;
            \item $\delta \gets \widetilde{b} - \mathbf{s}^{\top}\widetilde{\mathbf{a}}$;
            \item \textbf{return} $1$ iff $\|\delta\|_{\infty} > q/4$.
        \end{enumerate}
        \end{mdframed}
        \caption{Decryption in $\Pi_{\text{PKE}}^{\text{LWE}}$}
        \label{fig:lwe-dec}
    \end{minipage}
\end{figure}

Now we show that $\Pi_{\text{PKE}}^{\text{LWE}}$ is a correct and secure PKE scheme.

\begin{proposition}\label{prop:ind-cpa-security}
    $\Pi_{\text{PKE}}^{\text{LWE}}$ is correct and IND-CPA secure under DLWE
    assumption.
\end{proposition}

\textbf{Proof.} (\textit{Correctness}) Let us first show correctness. Assume $\mathsf{sk}=\mathbf{s}$
and $\mathsf{pk}=(A,\mathbf{b})$ is the valid key-pair (that is, $A\mathbf{s}+\mathbf{e}=\mathbf{b}$ for ``small'' $\mathbf{e} \leftarrowS \chi^m$) and we formed
$(\widetilde{\mathbf{a}}, \widetilde{b})$ according to \Cref{fig:lwe-enc}. Let 
us try to decode this cipher by means of \Cref{fig:lwe-dec}:
\begin{align*}
    \delta &= \mathbf{b}^{\top}\Delta\mathbf{r} + \Delta z' + m\lceil q/2 \rfloor - 
    \mathbf{s}^{\top}A^{\top}\Delta\mathbf{r} - \mathbf{s}^{\top}\Delta\mathbf{z} \\
    &= (\mathbf{b} - A\mathbf{s})^{\top}\Delta \mathbf{r} + (\Delta z' - \mathbf{s}^{\top}\Delta\mathbf{z}) + m\lceil q/2 \rfloor \\
    &= \textcolor{blue!70!black}{(\mathbf{e}^{\top}\Delta \mathbf{r} - \mathbf{s}^{\top}\Delta\mathbf{z} + \Delta z')} + m\lceil q/2 \rfloor
\end{align*}

Now note that the quantity marked in \textcolor{blue!70!black}{blue} is small.
Indeed, let us show why. Assume that $\chi$ produces elements of $\mathbb{Z}_q$
with the modulus not exceeding $\eta$. Then\footnote{In the sequel, we will show
that $\mathbf{s}$ can be taken from the same distribution as the error term
$\mathbf{e}$ (that is, from $\chi$), making upper bound even smaller.},
\begin{equation*}
    \|\textcolor{blue!70!black}{\mathbf{e}^{\top}\Delta \mathbf{r} - \mathbf{s}^{\top}\Delta\mathbf{z} + \Delta z'}\|_{\infty} \leq \|\mathbf{e}^{\top}\Delta \mathbf{r}\|_{\infty} + \|\mathbf{s}^{\top}\Delta\mathbf{z}\|_{\infty} + \|\Delta z'\|_{\infty} \leq n\eta^2 + n\eta + \eta \leq 2n\eta.
\end{equation*}

Clearly, $2n\eta$ is a small quantity (assuming $\eta \ll q/8n$). Now, if $m=0$, we
get that $\|\delta\|_{\infty} = \|(\textcolor{blue!70!black}{\text{blue part}})\|_{\infty} < q/4$, so
decryption works. If $m=1$, then $\delta = \lceil q/2\rfloor +
(\textcolor{blue!70!black}{\text{blue part}})$, which is a large value, 
so the decryption outputs $1$, as expected.

(\textit{IND-CPA Security}). Suppose $\mathcal{A}$ solves the
$\underline{\mathsf{Game} \; \mathsf{INDCPA}}_{\Pi}^{\mathcal{A}}$ with a
non-negligible probability. We construct the $\mathcal{B}^{\mathcal{A}}$ that
breaks the $\underline{\mathsf{Game} \;
\mathsf{DLWE}}_{\mathsf{LWEGen}}^{\mathcal{B}}$ using $\mathcal{A}$ as an
oracle. The high-level idea of the proof is the following: suppose $\mathcal{B}$
is given a tuple $(A,\mathbf{b})$. If $(A,\mathbf{b})$ is a valid LWE instance,
then it is a valid public key in $\Pi_{\text{PKE}}^{\text{LWE}}$, so
$\mathcal{B}$ might run $\mathcal{A}$ on $\mathsf{pk}^* \gets (A,\mathbf{b})$.
Then, if $\mathcal{A}$ correctly guesses the bit, then likely $\mathsf{pk}^*$
was a valid key. In turn, if $\mathsf{pk}^*$ is a ``garbage'', with probability
approximately $1/2$, the $\mathcal{A}$ would lose. Thus, $\mathcal{B}$ decision
will be based on whether $\mathcal{A}$ guesses the correct bit or not. We
encapsulate this in \Cref{fig:ind-cpa-to-dlwe}. One can then show:
\begin{equation*}
    \mathsf{Adv}^{\text{IND-CPA}}_{\mathcal{A},\Pi}(\lambda) \leq 2\mathsf{Adv}^{\text{DLWE}}_{\mathcal{B}^{\mathcal{A}},\mathsf{LWEGen}}(\lambda). 
    \quad \qedwhite
\end{equation*}

\begin{figure}[ht]
\centering
\begin{varwidth}{0.7\linewidth}
\begin{mdframed}[linecolor=black,
                 linewidth=1.0pt,
                 skipabove=10pt,
                 skipbelow=10pt,
                 leftmargin=0pt,
                 rightmargin=0pt]
\colorbox{purple!10}{\underline{$\mathcal{B}^{\mathcal{A}}_{\text{DLWE}}(1^{\lambda}, A, \mathbf{b}) \to \{0,1\}$:}}
\begin{enumerate}
    \item pass $\mathsf{pk}^* \gets (A,\mathbf{b})$ to $\mathcal{A}$;
    \item get $(m_0,m_1)$ from $\mathcal{A}$; 
    \item choose $d \leftarrowS \{0,1\}$;
    \item send $c_d \gets \mathsf{Enc}(\mathsf{pk}^*, m_d)$ to $\mathcal{A}$;
    \item $\mathcal{A}$ responds with $\widehat{d}$;
    \item \textbf{return} $d =_? \widehat{d}$. 
\end{enumerate}
\end{mdframed}
\end{varwidth}
\caption{Algorithm for finding the solution to DLWE instance given oracle to the IND-CPA.}
\label{fig:ind-cpa-to-dlwe}
\end{figure}

\subsection{Relation to Lattices}

The reasonable question that the reader might rightfully have at this point:
where are the lattices? So far, we have spent the whole \Cref{section:lattices}
describing lattices, while the current section does not explicitly use them in
any way.

In fact, turns out that the LWE problem can be viewed as the lattice problem!

\begin{definition}
    For the LWE instance $(A,\mathbf{b})$ over parameters $(n,m,q,\chi)$, we
    define the associated lattice $\Lambda_q(A)$ as follows:
    \begin{equation*}
        \Lambda_q(A) \triangleq \{\mathbf{y} \in \mathbb{Z}^m: \mathbf{y} = A\mathbf{x} \; (\text{mod}\; q) \; \text{for some} \; \mathbf{x} \in \mathbb{Z}^n\} = A\mathbb{Z}^n + q\mathbb{Z}^m.
    \end{equation*}
\end{definition}

Now note that the vector $\mathbf{b}$ does not lie on the lattice $\Lambda_q(A)$
(unless $\mathbf{e}=\mathbf{0}$ which makes LWE instance unsecure). But note
that due to the modular equation, $\mathbf{b} = A\mathbf{s} + q\mathbf{z} +
\mathbf{e}$ for some $\mathbf{z} \in \mathbb{Z}^m$ and since $A\mathbf{s} +
q\mathbf{z} \in \Lambda_q(A)$, the distance from $\mathbf{b}$ to $\Lambda_q(A)$
is bounded by $\|\mathbf{e}\|$. That said, although $\mathbf{b}$ is not in the
lattice, it is very close to it! Therefore, almost surely, the closest vector to
$\mathbf{b}$ is $A\mathbf{s} + q\mathbf{z}$. Now imagine that we know how to
solve CVP defined in \Cref{section:approximate-cvp}. Then, based on
$\mathbf{b}$, we can derive $A\mathbf{s}+q\mathbf{z}$, and then easily obtain
$\mathbf{s}$ by solving a system of linear equations. Note that the
approximation factor depends on the underlying distribution $\chi^m$, from which
$\mathbf{e}$ is sampled. 

This way, lattices are crucial for cryptanalysis of LWE-based schemes.

\subsection{Variations of LWE}

Let us analyze the efficiency of the PKE scheme provided in \Cref{def:lwe-pke}.
\begin{proposition}
    In the PKE scheme $\Pi_{\text{PKE}}^{\text{LWE}}$ over LWE instance
    $(n,m,q,\chi)$, the public key size is $O(m\log q)$, secret key size is
    $O(n\log q)$, and the codeword size is $O(m\log q)$: therefore, sending a
    $B$-bit message takes $O(B m \log q)$ space.
\end{proposition}

\begin{remark}\label{remark:seeded-a}
    One might ask why the public key size is not $O(nm\log q)$, given that we need
to store the matrix $A$. Fortunately, $A$ is sampled randomly, so instead of
storing the whole matrix, one might store the seed from which one can derive
$A$ instead.
\end{remark}

Unfortunately, such key and ciphertext sizes are significantly larger than those
of elliptic curve-based schemes (in fact, if $\lambda$ is taken as a security
parameter, then ciphertext size scales as $O(\lambda^2)$). This led to the
natural question: can we reduce these sizes?

\subsubsection{Ring-LWE}

Turns out that one possible answer to this question is to use objects other than
$\mathbb{Z}_q^n$ as elements that possess certain useful structure: for
instance, the popular choice is to use $\mathcal{R}_q :=
\mathcal{R}/q\mathcal{R}$ for $\mathcal{R} = \mathbb{Z}[T]/(T^n+1)$ and $n$
being a power-of-two. In such case, one might think of $\mathcal{R}$ either by: 
\begin{enumerate}
    \item[(a)] $(n-1)$-degree polynomials with operations modulo $T^n+1$;
    \item[(b)] the ring of integers $\mathbb{Z}(\zeta_{2n})$ since $\mathcal{R}
\cong \mathbb{Z}[T]/(\Phi_{2n}(T)) \cong \mathbb{Z}(\zeta_{2n})$ where $\zeta_N
:= e^{2\pi i/N}$.
\end{enumerate}

Now, instead of having $(A,\mathbf{b}) \in \mathbb{Z}_q^{m \times n} \times
\mathbb{Z}_q^m$ as a public key, we will store the public key as a pair $(a,b)
\in \mathcal{R}_q^2$ where $b=as+e$, but over the polynomial space. However, we
should clarify what ``small'' means in the context of $\mathcal{R}_q$ (since $e$
here comes from some small distribution). For this reason, given $f(T)=\sum_{i
\in (n)}f_iT^i \in \mathcal{R}_q$, define:
\begin{equation*}
    \|f\|_{\infty} \triangleq \max_{i \in (n)}\|f_i\|_{\infty}.
\end{equation*}
So in other words, polynomial $f(T) \in \mathcal{R}_q$ is large as long as its
largest coefficient is large. 

Further, we identify Ring-LWE instance by a tuple of $(\mathcal{R},n,q,\chi)$,
where (\textit{i}) additive group of $\mathcal{R}$ is isomorphic to
$\mathbb{Z}^n$, (\textit{ii}) instead of $\mathbb{Z}_q^n$ in ``plain'' LWE, we
use field $\mathcal{R}_q := \mathcal{R}/q\mathcal{R}$, and (\textit{iii}) $\chi$
is a (small) random distribution over $\mathcal{R}_q$.

\begin{proposition}
    Using Ring-LWE structure, we can reduce ciphertext size for a $n$-bit message
    of PKE Scheme from \Cref{def:lwe-pke} from $O(nm\log q)$ to $O(n\log q)$.
    Moreover, computing expression $as+e$ now takes only $O(n\log n)$ complexity
    instead of $O(nm)$.
\end{proposition}

The reason for $O(n\log n)$ complexity for computing $as+e$ comes from the fact
that we can utilize Number Theoretic Transform (NTT/FFT) for fast polynomial
multiplication. The reason for reduced ciphertext size is less evident, but the
idea is the following: instead of adding a $m\lceil q/2 \rfloor$ term to the
second ciphertext piece in \Cref{fig:lwe-enc}, we are going to represent $n$
bits of message $(m_i)_{i \in (n)}$ as a polynomial $m(T) := \sum_{i \in
(n)}m_iT^i \in \mathcal{R}_q$. We specify details in the definition below.

\begin{definition}\label{def:rlwe-pke}
    Based on Ring-LWE instance with parameters $(\mathcal{R}, n, q,\chi)$, define the PKE scheme
    $\Pi_{\text{PKE}}^{\text{RLWE}}$ over message space $\{0,1\}^n$ as follows: $\mathsf{KeyGen}$ is specified in
    \Cref{fig:ring-keygen}, $\mathsf{Enc}$ in \Cref{fig:ring-lwe-enc}, and $\mathsf{Dec}$
    in \Cref{fig:ring-lwe-dec}.
\end{definition}

\begin{figure}[ht]
\centering
\footnotesize 

\begin{minipage}[t]{0.32\linewidth}
    \begin{mdframed}[linecolor=black, linewidth=1.0pt,
                     skipabove=5pt, skipbelow=5pt,
                     innerleftmargin=5pt, innerrightmargin=5pt]
        \centering
        \colorbox{purple!10}{\underline{$\mathsf{KeyGen}(1^{\lambda}) \to (\mathsf{pk},\mathsf{sk})$:}}
        \begin{enumerate}
            \item $(s,a) \leftarrowS \mathcal{R}_q^2$;
            \item $e \leftarrowS \chi$ (over $\mathcal{R}_q$);
            \item $b \gets as + e \in \mathcal{R}_q$;
            \item \textbf{return} $((a,b),s)$.
        \end{enumerate}
    \end{mdframed}
    \caption{Ring-LWE-based Key Generation}
    \label{fig:ring-keygen}
\end{minipage}
\hfill
\begin{minipage}[t]{0.32\linewidth}
    \begin{mdframed}[linecolor=black, linewidth=1.0pt,
                     skipabove=5pt, skipbelow=5pt,
                     innerleftmargin=5pt, innerrightmargin=5pt]
        \centering
        \colorbox{green!10}{\underline{$\mathsf{Enc}(\mathsf{pk}, (m_i)_{i \in (n)} \in \{0,1\}^n) \to c$:}}
        \begin{enumerate}
            \item $m \gets \sum_{i \in (n)}m_iT^i \in \mathcal{R}_q$;
            \item $\Delta r, \Delta z, \Delta z' \leftarrowS \chi$;
            \item $\widetilde{a} \gets a\Delta r + \Delta z$;
            \item $\widetilde{b} \gets b\Delta r  + \Delta z' + \lceil \frac{q}{2} \rfloor m$;
            \item \textbf{return} $c=(\widetilde{a}, \widetilde{b})$.
        \end{enumerate}
    \end{mdframed}
    \caption{Ring-LWE-based Encryption} 
    \label{fig:ring-lwe-enc}
\end{minipage}
\hfill
\begin{minipage}[t]{0.32\linewidth}
    \begin{mdframed}[linecolor=black, linewidth=1.0pt,
                     skipabove=5pt, skipbelow=5pt,
                     innerleftmargin=5pt, innerrightmargin=5pt]
        \centering
        \colorbox{green!10}{\underline{$\mathsf{Dec}(\mathsf{sk}, c) \to m$:}}
        \begin{enumerate}
            \item Parse $(\widetilde{a}, \widetilde{b}) \gets c$;
            \item $\delta \gets \widetilde{b} - s\widetilde{a}$;
            \item \textbf{return} $1$ iff $\|\delta\|_{\infty} > \frac{q}{4}$.
        \end{enumerate}
    \end{mdframed}
    \caption{Ring-LWE-based Decryption}
    \label{fig:ring-lwe-dec}
\end{minipage}
\end{figure}

Another valid question is whether introducing such structure should worsen the
security: so in other words, given appropriate Ring-LWE instance
$(\mathcal{R},n,q,\chi)$, is $\Pi_{\text{PKE}}^{\text{RLWE}}$ secure? Turns out
that there are no algorithms for breaking analogous Ring-LWE-based assumption
significantly faster than ``plain'' version of the LWE. To put this assumption
formally, we slighly modify \Cref{def:search-lwe-assumption}.

\begin{definition}
    For the given adversary $\mathcal{A}$ and Ring-LWE parameters
    $(\mathcal{R},n,q,\chi)$, define $\underline{\mathsf{Game} \;
    \mathsf{SRLWE}}^{\mathcal{A}}$ as specified in~\Cref{fig:search-rlwe}. Define
    the advantage of $\mathcal{A}$ as
    \begin{equation*}
        \mathsf{Adv}^{\text{SRLWE}}_{\mathcal{A}}(\lambda) = \text{Pr}\left[\underline{\mathsf{Game} \;
    \mathsf{SRLWE}}^{\mathcal{A}}(\lambda) = 1\right].
    \end{equation*}

    We say that \textbf{Search-Ring-LWE assumption} (SRLWE) holds if
    $\mathsf{Adv}^{\text{SRLWE}}_{\mathcal{A}}(\lambda)=\text{negl}(\lambda)$
    for all PPT adversaries $\mathcal{A}$. Decisional version of Ring-LWE is
    defined similarly to \Cref{def:decision-lwe-assumption}.
\end{definition}

\begin{figure}[ht]
\centering
\begin{varwidth}{0.4\linewidth}
\begin{mdframed}[linecolor=black, linewidth=1pt,
                         skipabove=10pt, skipbelow=10pt]
\colorbox{green!10}{\underline{$\mathsf{Game} \; \mathsf{SRLWE}^{\mathcal{A}}$:}}
\begin{enumerate}
    \item $(s,a) \leftarrowS \mathcal{R}_q^2$;
    \item $e \leftarrowS \chi$;
    \item $b \gets as + e$;
    \item $\widehat{s} \gets \mathcal{A}(a,b)$;
    \item \textbf{return} $s =_? \widehat{s}$.
\end{enumerate}
\end{mdframed}
\end{varwidth}
\caption{Search-Ring-LWE Game} 
\label{fig:search-rlwe}
\end{figure}

We leave the following proposition without a proof as it closely resembles one
specified in Proposition~\ref{prop:ind-cpa-security}. For concrete details and
survey on possible attacks, see for example \cite{lac}.
\begin{proposition}
    $\Pi_{\text{PKE}}^{\text{RLWE}}$ is correct and IND-CPA secure under Decisional Ring-LWE
    assumption.
\end{proposition}

\subsubsection{Module-LWE}

The idea of Module-LWE is to in a way ``blend'' together Ring-LWE and plain-LWE.
Namely, instead of using $\mathcal{R}_q$ in Ring-LWE, we will use vectors from
$\mathcal{R}_q^{\ell}$. Note that equipping $\mathcal{R}_q$ with a vector
structure naturally induces metric notions we have used before: e.g., for
$\bm{f}=(f_1,\dots,f_{\ell}) \in \mathcal{R}_q^{\ell}$, we have
$\|\bm{f}\|_{\infty} := \max_{i \in [\ell]}\|f_i\|_{\infty}$ and $\langle
\bm{f}, \bm{g} \rangle = \sum_{i \in [\ell]}f_ig_i$ is a valid inner product
that produces output in $\mathcal{R}_q$. 

The analogy to plain-LWE equation $\mathbf{b} = A\mathbf{s}+\mathbf{e}$ for $A
\in \mathbb{Z}_q^{n \times m}$, $\mathbf{s} \in \mathbb{Z}_q^m$, and
$\mathbf{b},\mathbf{e} \in \mathbb{Z}_q^n$, is now $\mathbf{b}_{\mathcal{R}} =
A_{\mathcal{R}}\mathbf{s}_{\mathcal{R}} + \mathbf{e}_{\mathcal{R}}$ for
$A_{\mathcal{R}} \in \mathcal{R}_q^{k \times \ell}$, $\mathbf{s} \in
\mathcal{R}_q^{\ell}$, and $\mathbf{s}_{\mathcal{R}},\mathbf{e}_{\mathcal{R}}
\in \mathcal{R}_q^{k}$, where operations are done inside ring $\mathcal{R}_q$.
We call $(A,\mathbf{b})$ a Module-LWE instance with parameters
$(\mathcal{R},k,\ell,q,\chi)$ if it is an output of $\mathsf{MLWEGen}_{\mathcal{R},k,\ell,q,\chi}$
defined in \Cref{fig:mlwe-gen}. Now we define Module-LWE assumption in the 
definition below.

\begin{definition}
    For the given adversary $\mathcal{A}$ and Module-LWE parameters
    $(\mathcal{R},k,\ell,q,\chi)$, define $\underline{\mathsf{Game} \;
    \mathsf{SMLWE}}^{\mathcal{A}}$ as specified in~\Cref{fig:mlwe-game}. Define
    the advantage of $\mathcal{A}$ as
    \begin{equation*}
        \mathsf{Adv}^{\text{SMLWE}}_{\mathcal{A}}(\lambda) = \text{Pr}\left[\underline{\mathsf{Game} \;
    \mathsf{SMLWE}}^{\mathcal{A}}(\lambda) = 1\right].
    \end{equation*}

    We say that \textbf{Search-Module-LWE assumption} (SMLWE) holds if
    $\mathsf{Adv}^{\text{SMLWE}}_{\mathcal{A}}(\lambda)=\text{negl}(\lambda)$
    for all PPT adversaries $\mathcal{A}$. Decisional Module-LWE is defined
    similarly to \Cref{def:decision-lwe-assumption}.
\end{definition}

\begin{figure}[ht]
    \centering
    \begin{minipage}{0.45\linewidth}
        \begin{mdframed}[linecolor=black, linewidth=1pt,
                         skipabove=10pt, skipbelow=10pt]
        \colorbox{green!10}{\underline{$\mathsf{MLWEGen}_{\mathcal{R},k,\ell,q,\chi}(\mathbf{s}) \to \mathcal{R}_q^{k \times \ell}
            \times \mathcal{R}_q^{k}$:}}
        \begin{enumerate}
            \item $A \leftarrowS \mathcal{R}_q^{k \times \ell}$;
            \item $\mathbf{e} \leftarrowS \chi$ (over $\mathcal{R}_q^k$);
            \item $\mathbf{b} \gets A\mathbf{s} + \mathbf{e}$;
            \item \textbf{return} $(A,\mathbf{b})$.
        \end{enumerate}
        \end{mdframed}
        \caption{Modular-LWE $(\mathcal{R},k,\ell,q,\chi)$ instance generation} 
        \label{fig:mlwe-gen}
    \end{minipage}
    \hfill
    \begin{minipage}{0.45\linewidth}
        \begin{mdframed}[linecolor=black, linewidth=1pt,
                         skipabove=10pt, skipbelow=10pt]
        \colorbox{green!10}{\underline{$\mathsf{Game}\;\mathsf{SMLWE}_{\mathsf{MLWEGen}}^{\mathcal{A}}(1^{\lambda})$:}}
        \begin{enumerate}
            \item $\mathbf{s} \leftarrowS \mathcal{R}_q^{\ell}$;
            \item $(A,\mathbf{b}) \leftarrowS \mathsf{MLWEGen}(\mathbf{s})$;
            \item $\widehat{\mathbf{s}} \gets \mathcal{A}(A,\mathbf{b})$;
            \item \textbf{return} $\mathbf{s} =_? \widehat{\mathbf{s}}$.
        \end{enumerate}
        \end{mdframed}
        \caption{Search Modular-LWE Game}
        \label{fig:mlwe-game}
    \end{minipage}
\end{figure}

At first glance, it might seem that Module-LWE looks much less efficient than
Ring-LWE: for instance, why should representing $A$ as a matrix
$\mathcal{R}_q^{k \times \ell}$ be any more efficient than using a single
$\mathcal{R}_q$? Notice that to adjust the security of Ring-LWE, one must
regulate $n$ --- the power of cyclotomic polynomial. However, it must be a power
of two, which makes such parameter adjusting highly problematic. For Module-LWE,
one typically fixes $(q,n)$ (which fully determine $\mathcal{R}_q$-arithmetic
efficiency), and then changes $\ell$ to set target security level.

This way, key sizes (which are primarily determined by a storage of $\mathbf{b}$
by Remark~\ref{remark:seeded-a}) turn out to be more or less equal for both
Ring-LWE and Module-LWE:
\begin{itemize}
    \item \textit{Dilithium}~\cite{dilithium}, being a Module-LWE, uses
    parameters $q=2^{23}-2^{13}+1$, $n=256$, and $(k,\ell) \in
    \{(3,2),(4,3),(5,4),(6,5)\}$; If we take, say $(k,\ell)=(4,3)$, we need $kn
    = 1024$ elements of $\mathbb{Z}_q$ to store $\mathbf{b} \in
    \mathcal{R}_q^{k}$.
    \item \textit{NewHope}~\cite{new-hope}, being a Ring-LWE, uses
    $q=2^{13}+2^{12}+1$ and $n=1024$. Notice that although we have a single
    polynomial $b \in \mathcal{R}_q$, we still need $1024$ elements of
    $\mathbb{Z}_q$ for storage.
\end{itemize}

Turns out that implementing analogous PKE scheme to ones in \Cref{def:lwe-pke}
and \Cref{def:rlwe-pke} based on Module-LWE assumptions leads to the Kyber PKE
description! We specify the details in \Cref{section:kyber}.

\subsection*{Acknowledgements}

Some of the arguments and ideas used in these notes were particularly inspired
by awesome courses ``Cryptography 101 with Alfred Menezes''\footnote{See
\url{https://cryptography101.ca/}.} and lecture materials by François Charles on
lattices during ICMU Summer School on Number Theory.

\subsection*{References}
\footnotesize
\renewcommand{\refname}{\vspace{-1em}}
\bibliographystyle{alphaurl}
\bibliography{refs-lattices}

\end{document}